% Chapter 1 exercises. Source PDF pages 57--62 / printed pages 42--47.
\MathBookSourcePage{57}
\subsection*{习题}
\label{sec:ch01-exercises}
\setcounter{exerciseitem}{0}

\begin{MBExercise}
\label{ex:ch01-exercise-01}
设 $\Omega$ 有界且 $1\leq p\leq q\leq\infty$. 证明 $L^q(\Omega)\subset L^p(\Omega)$. (提示: 使用 Hölder 不等式.) 举例说明当 $p<q$ 时该包含为真包含, 而当 $\Omega$ 无界时该包含不成立.
\end{MBExercise}

\begin{MBExercise}
\label{ex:ch01-exercise-02}
证明区域 $\Omega$ 上的有界连续函数集合在范数 $\|\cdot\|_{L^\infty(\Omega)}$ 下构成 Banach 空间.
\end{MBExercise}

\begin{MBExercise}
\label{ex:ch01-exercise-03}
设 $\Omega$ 有界且 $f_j\to f$ 于 $L^p(\Omega)$. 使用 Hölder 不等式证明
\[
\MathBookFormulaID{formula-ch01-exercise-03-integral-limit}
\int_\Omega f_j(x)\,dx\longrightarrow\int_\Omega f(x)\,dx
\qquad\text{as }j\to\infty.
\]
\end{MBExercise}

\begin{MBExercise}
\label{ex:ch01-exercise-04}
这里补充 $n$ 维极坐标的知识. 归纳地定义从 $\Omega_k\subset\mathbb R^k$ 的子集到 $\mathbb R^{k+1}$ 单位球面的映射 $x^k$:
\[
\MathBookFormulaID{formula-ch01-exercise-04-spherical-recursion}
(x_1^k,\ldots,x_{k+1}^k)(\omega,\phi)
:=((\sin\phi)x^{k-1}(\omega),\cos\phi),
\]
其中
\[
\MathBookFormulaID{formula-ch01-exercise-04-angular-domains}
\Omega_k:=\{(\omega,\phi):\omega\in\Omega_{k-1},\ 0\leq\phi<\pi\},
\]
$x^1(\omega)=(\cos\omega,\sin\omega)$ 且 $\Omega_1=\{\omega:0\leq\omega<2\pi\}$. 对 $r\geq0$ 和 $\omega\in\Omega_{n-1}$, 定义极坐标映射 $X(\omega,r):=rx^{n-1}(\omega)$. 证明
\[
\MathBookFormulaID{formula-ch01-exercise-04-polar-jacobian}
\left|\det\frac{\partial X(\omega,r)}{\partial(\omega,r)}\right|
=r^{n-1}\prod_{j=2}^{n-1}(\sin\omega_j)^{j-1}
\]
对 $\omega=(\omega_1,\ldots,\omega_{n-1})\in\Omega_{n-1}$ 和 $r\geq0$ 成立. (提示: 先对 $n=2$ 计算, 再对 $k$ 归纳证明以下恒等式.)
\[
\MathBookFormulaID{formula-ch01-exercise-04-jacobian-recursion}
\det\begin{pmatrix}\dfrac{\partial x^{k+1}}{\partial(\omega,\phi)}\\[2pt]x^{k+1}(\omega,\phi)\end{pmatrix}
=(\sin\phi)^k
\det\begin{pmatrix}\dfrac{\partial x^k}{\partial\omega}\\[2pt]x^k(\omega)\end{pmatrix}.
\]
\end{MBExercise}

\begin{MBExercise}
\label{ex:ch01-exercise-05}
设 $n$ 为正整数, $\rho$ 为定义在 $0<r\leq1$ 上的非负光滑函数, 且满足
\[
\MathBookFormulaID{formula-ch01-exercise-05-radial-integrability}
\lim_{\epsilon\to0^+}\int_\epsilon^1\rho(r)r^{n-1}\,dr<\infty.
\]
在 $\Omega=\{x\in\mathbb R^n:|x|<1\}$ 上定义 $f(x)=\rho(|x|)$. 证明 $f\in L^1(\Omega)$. (提示: 使用单调收敛定理、极坐标(参见习题~\ref{ex:ch01-exercise-04})和 Fubini 定理.)
\end{MBExercise}

\begin{MBExercise}
\label{ex:ch01-exercise-06}
设 $\Omega=[0,1]$ 且 $1\leq p<\infty$. 证明 \eqref{eq:ch01-f-log-series} 定义的函数 $f$ 属于 $L^p(\Omega)$, 并且 $\|f-f_j\|_{L^p(\Omega)}\to0$, $j\to\infty$. (提示: 先证明 $\log x\in L^p(\Omega)$, 再使用 $L^p(\Omega)$ 为 Banach 空间这一事实.)
\end{MBExercise}

\MathBookSourcePage{58}
\begin{MBExercise}
\label{ex:ch01-exercise-07}
选取你喜欢的稠密序列 $\{r_n\}$. 对不同的 $j$ 画出 \eqref{eq:ch01-fj-log-series} 定义的函数 $f_j$ 的图像.
\end{MBExercise}

\begin{MBExercise}
\label{ex:ch01-exercise-08}
证明对所有 $x\neq0$ 和 $|\alpha|=1$,
\[
\MathBookFormulaID{formula-ch01-exercise-08-radial-derivative}
\left(\frac{\partial}{\partial x}\right)^\alpha|x|=\frac{x^\alpha}{|x|}.
\]
\end{MBExercise}

\begin{MBExercise}
\label{ex:ch01-exercise-09}
证明命题~\ref{prop:ch01-classical-equals-weak}. (提示: 使用分部积分和 $C^0(\Omega)\subset L_{\mathrm{loc}}^1(\Omega)$ 这一事实.)
\end{MBExercise}

\begin{MBExercise}
\label{ex:ch01-exercise-10}
证明例~\ref{ex:ch01-absolute-value-weak-derivative} 中函数 $f$ 的二阶及更高阶弱导数不存在.
\end{MBExercise}

\begin{MBExercise}
\label{ex:ch01-exercise-11}
证明例~\ref{ex:ch01-radial-weak-derivative} 中的条件蕴含
\[
\MathBookFormulaID{formula-ch01-exercise-11-rho-limit}
\lim_{r\to0}r^{n-1}|\rho(r)|=0.
\]
(提示: 先证明 $r^{n-1}|\rho(r)|\leq\int_r^1t^{n-1}|\rho'(t)|\,dt+C\leq\widetilde C$. 若极限不为零, 则存在趋于零的点 $r_j$, 使对适当的 $\gamma>0$,
\[
\MathBookFormulaID{formula-ch01-exercise-11-contradiction-bound}
\begin{aligned}
\left|r_j^{n-1}\int_{r_j}^{r_j+\gamma r_j}\rho'(t)\,dt\right|
&=\left|r_j^{n-1}\bigl(\rho(r_j+\gamma r_j)-\rho(r_j)\bigr)\right|\\
&\geq\left|r_j^{n-1}
\bigl(C(r_j+\gamma r_j)^{1-n}-cr_j^{1-n}\bigr)\right|=\delta>0.
\end{aligned}
\]
说明这与条件矛盾. 或者应用 Hardy 不等式 (Stein 1970).)
\end{MBExercise}

\begin{MBExercise}
\label{ex:ch01-exercise-12}
验证例~\ref{ex:ch01-radial-weak-derivative} 中所断言的弱导数存在性. (提示: 使用习题~\ref{ex:ch01-exercise-11}, 习题~\ref{ex:ch01-exercise-08}, 习题~\ref{ex:ch01-exercise-05}, 散度定理以及 Lebesgue 积分的一个极限定理.)
\end{MBExercise}

\begin{MBExercise}
\label{ex:ch01-exercise-13}
对给定实数 $r$, 令 $f(x)=|x|^r$. 证明当 $r>1-n$ 时, $f$ 在单位球上有一阶弱导数.
\end{MBExercise}

\begin{MBExercise}
\label{ex:ch01-exercise-14}
证明 $\operatorname{Lip}(\Omega)\subset W_\infty^1(\Omega)$. (提示: $f\in\operatorname{Lip}(\Omega)$ 对每个变量分别而言当然也是 Lipschitz 连续的, 因而几乎处处有偏导数. 使用 Fubini 定理和反证法说明这些偏导数在 $\Omega$ 上必本质有界. 利用 Lipschitz 函数绝对连续且绝对连续函数可以分部积分, 说明这些导数实际上是弱导数.) 对这一包含关系, Lipschitz 范数与 Sobolev $W_\infty^1(\Omega)$ 范数之间等价关系的常数是多少?
\end{MBExercise}

\begin{MBExercise}
\label{ex:ch01-exercise-15}
设 $\Omega$ 为凸集. 证明 $W_\infty^1(\Omega)\subset\operatorname{Lip}(\Omega)$. (提示: 给定 $f\in W_\infty^1(\Omega)$, $\phi\in\mathcal D(\Omega)$ 和 $y,z\in\Omega$, 写成
\[
\MathBookFormulaID{formula-ch01-exercise-15-directional-identity}
f(y)-f(z)=\lim_{\epsilon\to0}\int_\Omega f(x)D\phi^\epsilon(x)\,dx,
\]
其中 $D$ 表示沿 $n_{y,z}=(y-z)/|y-z|$ 方向的方向导数, 且
\[
\MathBookFormulaID{formula-ch01-exercise-15-phi-epsilon}
\phi^\epsilon(x)=\epsilon^{-n}\int_0^\infty
\phi\left(\frac{x-y-tn_{y,z}}\epsilon\right)
\phi\left(\frac{x-z-tn_{y,z}}\epsilon\right)\,dt.
\]
注意: 必须证明上述恒等式, 参见注~\ref{rem:ch01-mollification-density} 中的光滑化论证和习题~\ref{ex:ch01-exercise-18}. 通过验证并使用恒等式
\[
\MathBookFormulaID{formula-ch01-exercise-15-phi-alternative}
\phi^\epsilon(x)=-\epsilon^{-n}|y-z|
\int_{-1}^0\phi\left(\frac{x-y-\tau(y-z)}\epsilon\right)\,d\tau,
\]
证明 $\phi^\epsilon\in\mathcal D(\Omega)$ 且 $\|\phi^\epsilon\|_{L^1(\mathbb R^n)}=|y-z|$. 应用弱导数的定义和 Hölder 不等式.)
\end{MBExercise}

\MathBookSourcePage{59}
\begin{MBExercise}
\label{ex:ch01-exercise-16}
设 $n=1$, $\Omega=[a,b]$, $f\in W_1^1(\Omega)$. 在假设 $f$ 在 $a,b$ 处连续的条件下, 证明
\[
\MathBookFormulaID{formula-ch01-exercise-16-fundamental-theorem}
\int_a^bD_w^1f(x)\,dx=f(b)-f(a).
\]
(提示: 使用定义弱导数的``分部积分''公式, 并选取适当序列 $\phi_j\in\mathcal D(\Omega)$, 使 $\phi_j\to1$ 于 $\Omega$, 参见习题~\ref{ex:ch01-exercise-15}.)
\end{MBExercise}

\begin{MBExercise}
\label{ex:ch01-exercise-17}
证明区间 $[a,b]$ 上的绝对连续函数属于 $W_1^1([a,b])$. (提示: 使用绝对连续函数可以分部积分这一事实.)
\end{MBExercise}

\begin{MBExercise}
\label{ex:ch01-exercise-18}
验证注~\ref{rem:ch01-mollification-density} 中的论断. (提示: 对所有 $\alpha$ 证明 $D^\alpha f_\epsilon=\epsilon^{-|\alpha|}f*(D^\alpha\phi)_\epsilon$; 证明 $f_\epsilon\to f$ 于 $L^p(\Omega)$; 并在适当条件下证明 $D^\alpha f_\epsilon=(D_w^\alpha f)*\phi_\epsilon=(D_w^\alpha f)_\epsilon$.)
\end{MBExercise}

\begin{MBExercise}
\label{ex:ch01-exercise-19}
若对某个点 $x$, 每个 $y\in\Omega$ 与 $x$ 的连线段都位于 $\Omega$ 内, 则称区域 $\Omega$ 关于 $x$ 为星形. 这里指闭线段, 因而必有 $x\in\Omega$. 在 $\Omega$ 为星形的假设下, 证明当 $1\leq p<\infty$ 时 $C^\infty(\overline\Omega)$ 在 $W_p^k(\Omega)$ 中稠密. (提示: 见 (Dupont and Scott 1979). 假设 $x=0$, 证明当 $\rho<1$ 时定义的``伸缩'' $u_\rho(y)=u(\rho y)$ 满足 $u_\rho\to u$ 于 $W_p^k(\Omega)$, $\rho\to1$. 像注~\ref{rem:ch01-mollification-density} 中那样光滑化 $u_\rho$, 得光滑函数 $u_{\rho,\epsilon}\to u$, $\rho\to1$, $\epsilon\to0$.)
\end{MBExercise}

\begin{MBExercise}
\label{ex:ch01-exercise-20}
证明 $n=1$ 时的 Sobolev 不等式. 也就是说, 设 $\Omega=[a,b]$, 证明
\[
\MathBookFormulaID{formula-ch01-exercise-20-one-dimensional-sobolev}
\|u\|_{L^\infty(\Omega)}\leq C\|u\|_{W_1^1(\Omega)}.
\]
(提示: 使用微积分基本定理和定理~\ref{thm:ch01-meyers-serrin-density}. 技术细节可参考迹不等式 \eqref{eq:ch01-trace-inequality-unit-disk} 的证明.) 常数 $C$ 如何依赖于 $a,b$?
\end{MBExercise}

\MathBookSourcePage{60}
\begin{MBExercise}
\label{ex:ch01-exercise-21}
设 $\Omega=[a,b]$(这里 $n=1$). 证明 $W_1^1(\Omega)$ 中所有函数都连续(即具有连续代表). (提示: 使用习题~\ref{ex:ch01-exercise-20} 和定理~\ref{thm:ch01-meyers-serrin-density}.)
\end{MBExercise}

\begin{MBExercise}
\label{ex:ch01-exercise-22}
证明 $W_1^1([a,b])$ 中的函数在 $[a,b]$ 上绝对连续. (提示: 使用习题~\ref{ex:ch01-exercise-16}、习题~\ref{ex:ch01-exercise-21} 以及 Lebesgue 积分的``连续性''.)
\end{MBExercise}

\begin{MBExercise}
\label{ex:ch01-exercise-23}
使用习题~\ref{ex:ch00-classical-existence} 中的提示, 证明给定 $f\in L^2(\Omega)$, 存在满足
\eqref{eq:ch00-boundary-value-problem} 的解 $u\in W_2^2(\Omega)$. 验证所有边界条件都有意义且得到满足, 并说明应在何种意义下解释该微分方程.
\end{MBExercise}

\begin{MBExercise}
\label{ex:ch01-exercise-24}
设 $u$ 是由习题~\ref{ex:ch00-classical-existence} 或习题~\ref{ex:ch01-exercise-23} 给出的
\eqref{eq:ch00-boundary-value-problem} 的解. 严格证明
\eqref{eq:ch00-weak-problem} 成立.
\end{MBExercise}

\begin{MBExercise}
\label{ex:ch01-exercise-25}
证明任意 Lipschitz 区域都满足注~\ref{rem:ch01-boundary-smooth-density} 中的线段条件.
\end{MBExercise}

\begin{MBExercise}
\label{ex:ch01-exercise-26}
令
\[
\MathBookFormulaID{formula-ch01-exercise-26-slit-domain}
\Omega=\{(x,y)\in\mathbb R^2:|x|<1,\ |y|<1,\ y\neq0
\text{ for }x\geq0\}.
\]
证明 $W_p^1(\Omega)$ 中存在不能由 $C^0(\overline\Omega)$ 中函数的极限得到的函数.
\end{MBExercise}

\begin{MBExercise}
\label{ex:ch01-exercise-27}
``尖点''区域 $\Omega=\{(x,y):0<x<1,\ 0<y<x^r\}$, $r>1$, 是否满足线段条件? 它的补集是否满足?
\end{MBExercise}

\begin{MBExercise}
\label{ex:ch01-exercise-28}
设 $\Omega$ 为平面上的单位圆盘. 证明存在 $u\in W_2^1(\Omega)$, 它在 $\partial\Omega$ 的一个稠密点集处为无穷大. (提示: 参见
\eqref{eq:ch01-f-log-series} 和例~\ref{ex:ch01-unbounded-sobolev-function}.)
\end{MBExercise}

\begin{MBExercise}
\label{ex:ch01-exercise-29}
设 $\Omega$ 为平面上的区域. 证明存在 $u\in W_2^1(\Omega)$, 它在 $\Omega$ 的一个稠密点集处为无穷大. (提示: 参见
\eqref{eq:ch01-f-log-series} 和例~\ref{ex:ch01-unbounded-sobolev-function}.)
\end{MBExercise}

\begin{MBExercise}
\label{ex:ch01-exercise-30}
设 $\Omega$ 如命题~\ref{prop:ch01-unit-disk-trace}, $p$ 为满足 $1\leq p\leq\infty$ 的实数. 证明存在常数 $C$, 使
\[
\MathBookFormulaID{formula-ch01-exercise-30-general-p-trace}
\|v\|_{L^p(\partial\Omega)}
\leq C\|v\|_{L^p(\Omega)}^{1-1/p}
\|v\|_{W_p^1(\Omega)}^{1/p}
\qquad\forall v\in W_p^1(\Omega).
\]
说明当 $p=\infty$ 时该式的含义.
\end{MBExercise}

\begin{MBExercise}
\label{ex:ch01-exercise-31}
设 $\Omega$ 为上半平面. 证明不存在形如
\[
\MathBookFormulaID{formula-ch01-exercise-31-impossible-scaling}
\|v\|_{L^2(\partial\Omega)}
\leq C\|v\|_{L^2(\Omega)}^\lambda
\|v\|_{W_2^1(\Omega)}^\mu
\]
的不等式, 除非 $\lambda\leq1/2$ 且 $\lambda+\mu=1$. (提示: 假设该不等式成立, 考虑函数 $u(x)=Uv(Lx)$, 其中 $U,L$ 任意, 从而得到矛盾.) 证明要在 $\lambda=\mu=1/2$ 时建立该不等式, 只须假设 $\|v\|_{L^2(\Omega)}=\|v\|_{W_2^1(\Omega)}=1$. 也就是说, 从这个特殊情形推出一般情形. (提示: 使用相同的缩放.) 用同样的思想证明
\[
\MathBookFormulaID{formula-ch01-exercise-31-seminorm-trace}
\|v\|_{L^2(\partial\Omega)}
\leq C\|v\|_{L^2(\Omega)}^{1/2}|v|_{W_2^1(\Omega)}^{1/2}.
\]
\end{MBExercise}

\MathBookSourcePage{61}
\begin{MBExercise}
\label{ex:ch01-exercise-32}
设 $a,b\in\mathbb R$. 证明对任意 $\epsilon>0$,
\[
\MathBookFormulaID{formula-ch01-exercise-32-young}
ab\leq\frac\epsilon2a^2+\frac1{2\epsilon}b^2.
\]
(提示: 展开 $(\sqrt\epsilon a-b/\sqrt\epsilon)^2\geq0$.) 应用该式证明
\[
\MathBookFormulaID{formula-ch01-exercise-32-two-term-bound}
a+b\leq\bigl((1+\epsilon)a^2+(1+1/\epsilon)b^2\bigr)^{1/2}.
\]
\end{MBExercise}

\begin{MBExercise}
\label{ex:ch01-exercise-33}
证明线性空间上的线性泛函集合本身是线性空间, 其中加法定义为 $(L_1+L_2)(v):=L_1(v)+L_2(v)$, 数乘定义为 $(aL)(v):=aL(v)$.
\end{MBExercise}

\begin{MBExercise}
\label{ex:ch01-exercise-34}
证明
\eqref{eq:ch01-dual-norm} 中的表达式在 Banach 空间的对偶空间上定义范数, 即验证定义~\ref{def:ch01-norm} 中的各条件.
\end{MBExercise}

\begin{MBExercise}
\label{ex:ch01-exercise-35}
设 $B,C$ 为两个 Banach 空间, $B\subset C$, 且该包含映射连续. 证明 $C'\subset B'$, 且该包含映射连续.
\end{MBExercise}

\begin{MBExercise}
\label{ex:ch01-exercise-36}
设 $\Omega$ 为开集且 $y\in\Omega$. 证明不存在 $f\in L_{\mathrm{loc}}^1(\Omega)$, 使
\[
\MathBookFormulaID{formula-ch01-exercise-36-no-dirac-density}
\phi(y)=\int_\Omega f(x)\phi(x)\,dx
\qquad\forall\phi\in\mathcal D(\Omega).
\]
\end{MBExercise}

\begin{MBExercise}
\label{ex:ch01-exercise-37}
设 $\Omega=[0,1]$. 证明
\[
\MathBookFormulaID{formula-ch01-exercise-37-dense-zero-value-smooth}
\{v\in C^1([0,1]):v(0)=0\}
\]
在 $\{v\in W_2^1(\Omega):v(0)=0\}$ 中稠密. (提示: 使用注~\ref{rem:ch01-boundary-smooth-density} 中的稠密性结果和 Sobolev 不等式.)
\end{MBExercise}

\begin{MBExercise}
\label{ex:ch01-exercise-38}
设 $\Omega=[0,1]$. 证明 $\{v\in C^1([0,1]):v'(1)=0\}$ 在 $W_2^1(\Omega)$ 中稠密. (提示: 使用注~\ref{rem:ch01-boundary-smooth-density} 中的稠密性结果, 取 $C^1([0,1])$ 中收敛到 $v$ 于 $W_2^1(\Omega)$ 的序列 $v_n$, 再在 $1$ 附近修改每个 $v_n$ 使其导数为零.)
\end{MBExercise}

\begin{MBExercise}
\label{ex:ch01-exercise-39}
设 $\Omega$ 为开集. 假设 $f\in L_{\mathrm{loc}}^1(\Omega)$ 且
\[
\MathBookFormulaID{formula-ch01-exercise-39-zero-distribution}
\int_\Omega f\phi\,dx=0
\qquad\forall\phi\in\mathcal D(\Omega).
\]
证明 $f=0$ 几乎处处成立. (提示: 用反证法. 否则, 说明存在 $\epsilon>0$, 使 $A=\{x\in\Omega:\pm f(x)>\epsilon\}$ 具有正测度. 用有限个球的并 $A_N$ 在测度意义下逼近 $A$, 并令 $\widetilde A_N$ 表示稍小的同心球的并, 它也在测度意义下逼近 $A$. 选取 $\phi_N\in\mathcal D(A_N)$, 使处处有 $0\leq\phi_N\leq1$, 并在 $\widetilde A_N$ 上有 $\phi_N\geq1/2$, 再对 $f\phi_N$ 在 $\Omega$ 上积分以得到矛盾. 如下构造 $\phi_N$: 写成 $A_N=\bigcup_jB_j$, $\widetilde A_N=\bigcup_j\widetilde B_j$, 并选取 $\psi_j\in\mathcal D(B_j)$, 使处处有 $0\leq\psi_j\leq1$, 且在 $\widetilde B_j$ 上 $\psi_j=1$. 定义
\[
\MathBookFormulaID{formula-ch01-exercise-39-partition-function}
\phi_N(x)=\sum_j\psi_j(x)\left(1+\sum_j\psi_j(x)\right)^{-1}.
\]
这可以看作一个近似单位分解论证.)
\end{MBExercise}

\MathBookSourcePage{62}
\begin{MBExercise}
\label{ex:ch01-exercise-40}
证明弱导数是唯一的. (提示: 使用习题~\ref{ex:ch01-exercise-39}.)
\end{MBExercise}

\begin{MBExercise}
\label{ex:ch01-exercise-41}
与习题~\ref{ex:ch01-exercise-39} 一样, 设 $\widetilde B_j\subset B_j$ 为同心球, $\psi_j\in\mathcal D(B_j)$, 处处有 $0\leq\psi_j\leq1$, 且在 $\widetilde B_j$ 上 $\psi_j=1$. 设 $\epsilon>0$, 定义
\[
\MathBookFormulaID{formula-ch01-exercise-41-partition-function}
\phi^\epsilon(x)=\sum_j\psi_j(x)
\left(\frac\epsilon{1-\epsilon}+\sum_j\psi_j(x)\right)^{-1}.
\]
证明处处有 $0\leq\phi^\epsilon\leq1$, 且在 $\bigcup_j\widetilde B_j$ 上 $\phi^\epsilon\geq1-\epsilon$. 能否构造 $\phi\in\mathcal D(\bigcup_jB_j)$, 使处处有 $0\leq\phi\leq1$, 且在 $\bigcup_j\widetilde B_j$ 上 $\phi=1$? 若有限并改为无限并, 可以得到什么结论?
\end{MBExercise}

\begin{MBExercise}
\label{ex:ch01-exercise-42}
令
\[
\MathBookFormulaID{formula-ch01-exercise-42-corner-function-domain}
v_\beta(r,\theta)=r^\beta\sin(\beta\theta),
\qquad
\Omega_\beta=\{(r,\theta):r<1,\ 0<\theta<\pi/\beta\},
\]
其中 $1/2\leq\beta<\infty$. 作为 $\beta$ 的函数, 确定使 $v_\beta\in W_p^k(\Omega_\beta)$ 的最优 $k,p$ 值. $\beta=1/2$ 的情形称为狭缝区域, $\beta=2/3$ 的情形称为凹角. $\beta=1$ 对应什么情形?
\end{MBExercise}

\begin{MBExercise}
\label{ex:ch01-exercise-43}
冰淇淋所用的有限圆锥 $C$ 可以表示为集合 $t\xi+B_{\alpha t}$ 在 $0<t<h$ 时的并, 其中 $\xi$ 是圆锥轴, $\alpha$ 度量夹角, $h$ 是高度, $B_t$ 是半径为 $t$ 的球. Lipschitz 区域 $\Omega$ 满足圆锥性质: 存在有限个有限圆锥 $C_i$, 使对所有 $x\in\Omega$, 都存在某个 $i$ 使 $x+C_i\subset\Omega$. 对这样的区域, 证明光滑函数在 $W_p^k(\Omega)$ 中稠密. (提示: 定义 $x$ 处的光滑化值时, 使用光滑化核在集合 $x+t\xi+B_{\alpha t}$ 上取平均, 即
\[
\MathBookFormulaID{formula-ch01-exercise-43-cone-mollification}
v_t(x):=\int_{B_{\alpha t}}v(x+y+t\xi)\rho_t(y)\,dy.
\]
证明 $v_t$ 趋于 $v$ 于 $W_p^k(\Omega)$.)
\end{MBExercise}

\begin{MBExercise}
\label{ex:ch01-exercise-44}
Cantor 函数可以由连续分片线性函数的极限序列定义. 设 $0<\epsilon<1$, 从 $f_0(x)=x$ 开始. 接着令 $f_1(x)$ 在以 $x=1/2$ 为中心、宽度为 $\epsilon$ 的区间内等于 $1/2$, 并在区间 $[0,1]$ 的两端等于 $f_0$. 已给定 $f_i$, 归纳定义 $f_{i+1}$. 在 $f_i$ 不为常数的每个区间中, 像处理原区间一样划分: 令 $f_{i+1}$ 在中心长度为 $2^{-i}\epsilon$ 的子区间上为常数, 其余部分用线性函数连接. 该常数取为中间值. 证明序列 $f_i$ 在函数空间 $\operatorname{Lip}^\alpha$ 中有界, 该空间中的函数对某个 $1>\alpha>0$ 满足一致估计
\[
\MathBookFormulaID{formula-ch01-exercise-44-holder-bound}
|f(x)-f(y)|\leq C|x-y|^\alpha.
\]
证明当 $\epsilon$ 趋于零时, 最优的 $\alpha$ 可以趋于一.
\end{MBExercise}
